\chapter{Análisis de los resultados de la encuesta}

Con el objetivo de relevar la percepción de los usuarios respecto de la
administración y optimización de recursos en entornos de computación en
la nube, se realizó una encuesta mediante Google Forms. Se obtuvieron un
total de 48 respuestas. El cuestionario permitió recopilar información
tanto de usuarios con experiencia práctica en la administración de
infraestructura cloud como de personas sin experiencia directa,
diferenciando posteriormente ambos grupos para evitar que la falta de
experiencia condicionara el análisis de las preguntas de carácter
técnico. Los resultados completos y las respuestas individuales se
presentan en el Anexo correspondiente.

En primer lugar, se analizó la existencia de dificultades relacionadas
con la utilización eficiente de los recursos cloud. De los 48
participantes, 15 manifestaron contar con experiencia práctica en la
administración de infraestructura en la nube. Dentro de este grupo, 10
participantes indicaron haber experimentado gastos inesperados o
sobrecostos asociados al uso de servicios cloud, mientras que los 5
restantes indicaron no haber atravesado esta situación. En consecuencia,
aproximadamente el 66,7 \% de los participantes con experiencia práctica
manifestó haber experimentado algún grado de sobrecosto. Este resultado
permite observar que los costos derivados de una utilización ineficiente
de los recursos constituyen una problemática presente entre los usuarios
con experiencia en este tipo de infraestructura.

Se consultó a los participantes con experiencia sobre la dificultad para
identificar recursos ineficientes dentro de una infraestructura cloud,
utilizando una escala de 1 a 5. Las respuestas se concentraron
principalmente en los valores intermedios y altos, obteniéndose un
promedio de 3,56 sobre 5. Esto permite caracterizar la identificación de
recursos ineficientes como una tarea de dificultad moderada a elevada.

En relación con las métricas utilizadas para evaluar el estado de los
recursos, los costos fueron el indicador mencionado con mayor frecuencia
entre los participantes con experiencia, seguido por el uso de CPU y el
almacenamiento. La memoria RAM también presentó una frecuencia elevada,
mientras que variables como la disponibilidad, el tráfico de red y el
uso de contenedores tuvieron una presencia menor. Estos resultados
respaldan la selección inicial de métricas consideradas por la solución
propuesta, particularmente aquellas relacionadas con el consumo de CPU,
memoria y almacenamiento, complementadas por información económica
asociada a los recursos.

Por otra parte, se relevó qué herramientas utilizan actualmente los
participantes con experiencia práctica para monitorear, administrar y
optimizar sus infraestructuras cloud. Entre las soluciones mencionadas
se encontraron herramientas ampliamente reconocidas en el mercado, como
Grafana, Amazon CloudWatch, Datadog, Prometheus, Kubernetes Dashboard,
New Relic y Terraform. Resulta particularmente relevante que varios de
los participantes con experiencia práctica indicaron utilizar
herramientas que constituyen alternativas o competidores directos
respecto de distintas funcionalidades contempladas en la solución
propuesta. Por ejemplo, Datadog y New Relic ofrecen capacidades de
monitoreo y observabilidad, mientras que Amazon CloudWatch proporciona
servicios de monitoreo sobre infraestructura de AWS. La presencia de
estas herramientas entre las respuestas permite considerar que una parte
de la muestra posee conocimiento y experiencia real con soluciones
existentes en el mercado, por lo que sus respuestas aportan una
perspectiva relevante para evaluar las necesidades y limitaciones de las
alternativas actuales. A su vez, la diversidad de herramientas
utilizadas evidencia que los usuarios recurren a diferentes soluciones
para cubrir necesidades específicas de monitoreo, visualización,
administración e infraestructura como código. Esta posible fragmentación
resulta relevante para la propuesta desarrollada, que busca integrar en
una misma plataforma el análisis de utilización de recursos, la
detección de ineficiencias, la generación de recomendaciones de
optimización y la automatización mediante infraestructura como código.

Respecto de las principales limitaciones percibidas en las herramientas
utilizadas actualmente, la complejidad fue el aspecto mencionado con
mayor frecuencia, seguida por los costos elevados y la poca capacidad de
visualización. También se identificaron dificultades relacionadas con la
integración entre herramientas y con la utilidad de algunas de las
recomendaciones proporcionadas por las soluciones existentes. Estos
resultados permiten observar que la problemática no se limita
exclusivamente a la disponibilidad de métricas, sino que también
comprende la capacidad de interpretar dicha información y transformarla
en acciones concretas de optimización. En este sentido, una solución que
únicamente presente métricas o detecte anomalías podría resultar
insuficiente si no proporciona información que facilite la toma de
decisiones.

La encuesta también permitió analizar el interés de los participantes
respecto de las funcionalidades que podría ofrecer una plataforma
orientada a la optimización de recursos cloud. Las funcionalidades
seleccionadas con mayor frecuencia fueron el monitoreo del uso de
recursos y la generación de recomendaciones para optimizar costos, ambas
mencionadas por 27 participantes. A continuación, se ubicó la detección
automática de desperdicio de recursos, seleccionada por 20
participantes. También se identificó interés en funcionalidades
relacionadas con la generación de reportes y visualizaciones, la
automatización de despliegues y las alertas. Estos resultados presentan
una correspondencia directa con los objetivos de la solución propuesta,
dado que las principales funcionalidades demandadas pueden organizarse
en un flujo compuesto por la observación del comportamiento de los
recursos, la detección de situaciones ineficientes y la generación de
recomendaciones orientadas a su optimización.

Un aspecto particularmente relevante fue la valoración general de una
plataforma destinada a administrar recursos cloud de manera más
eficiente. Sobre el total de 48 participantes, 37 calificaron su
utilidad con un valor de 4 o 5 sobre 5, lo que representa
aproximadamente el 77,1 \% de la muestra, mientras que la valoración
promedio obtenida fue de 4 sobre 5. Entre los 15 participantes que
declararon contar con experiencia práctica en la administración de
infraestructura cloud, 12 calificaron la utilidad de la plataforma con 4
o 5, lo que representa el 80 \% de este grupo. Por otro lado, las cinco
valoraciones más bajas, correspondientes a valores de 1 o 2, provinieron
tanto de participantes con experiencia como de participantes sin
experiencia, por lo que no se observa que el rechazo hacia la propuesta
esté concentrado exclusivamente en alguno de estos grupos. En conjunto,
estos resultados evidencian un nivel elevado de interés potencial en una
solución orientada a facilitar la administración eficiente de recursos
cloud.

Sin embargo, la aceptación de las recomendaciones automatizadas presentó
un comportamiento más moderado. La confianza promedio declarada por los
participantes fue de 3,44 sobre 5, mientras que el 50 \% de la muestra
otorgó una valoración de 4 o 5. Al analizar específicamente a los
participantes con experiencia práctica en infraestructura cloud, este
porcentaje se incrementó al 66,7 \%. Esto permite identificar una
diferencia entre el interés por utilizar una herramienta de optimización
y la disposición a delegar decisiones de infraestructura en un sistema
automatizado. En consecuencia, los resultados sugieren que la
automatización constituye una funcionalidad valorada, pero debe estar
acompañada por mecanismos que permitan al usuario comprender, revisar y
controlar las acciones recomendadas.

La seguridad se identificó como el principal factor que condiciona la
confianza en las recomendaciones automatizadas, siendo mencionada por 24
participantes. Otros factores señalados fueron el riesgo de errores, la
experiencia previa, los costos y la falta de transparencia. Este
resultado tiene implicancias directas sobre el diseño de la solución
propuesta. En lugar de plantear una modificación automática e
irreversible de la infraestructura, resulta conveniente adoptar un
enfoque basado en recomendaciones explicables y controladas por el
usuario. Bajo este enfoque, el sistema puede detectar una situación
potencialmente ineficiente, explicar los datos que justifican la
recomendación, proponer una acción de optimización y, posteriormente,
permitir al usuario decidir si desea aplicarla.

Esta consideración resulta especialmente relevante para la integración
con Terraform planteada en el proyecto. En lugar de modificar
directamente los recursos productivos, la plataforma puede utilizar los
resultados del análisis para generar una propuesta de infraestructura
como código optimizada. De esta manera, el usuario conserva el control
sobre el cambio y puede revisar el código generado antes de aplicarlo.
Este mecanismo permite combinar automatización con supervisión humana,
reduciendo el riesgo asociado a modificaciones automáticas sobre
infraestructura crítica.

En conjunto, los resultados obtenidos permiten identificar tres aspectos
principales. En primer lugar, existe evidencia de que los usuarios con
experiencia en infraestructura cloud enfrentan problemas relacionados
con sobrecostos y con la identificación de recursos ineficientes. En
segundo lugar, las funcionalidades consideradas más relevantes por los
participantes monitoreo, optimización de costos y detección automática
de desperdicio presentan una correspondencia directa con los objetivos
de la solución propuesta. Finalmente, aunque existe una valoración
positiva hacia una plataforma de este tipo, la confianza en las
recomendaciones automatizadas se encuentra condicionada principalmente
por cuestiones de seguridad y transparencia.

Por lo tanto, los resultados de la encuesta proporcionan evidencia
favorable para continuar con el desarrollo de una plataforma orientada a
la detección de ineficiencias en infraestructuras cloud y a la
generación de recomendaciones de optimización. Al mismo tiempo, permiten
establecer criterios de diseño que trascienden la detección de
anomalías: las recomendaciones deben ser comprensibles, justificables y
revisables por el usuario, y las acciones de modificación de
infraestructura deberían mantener un mecanismo de supervisión antes de
su aplicación. De esta manera, la encuesta no solo contribuye a validar
la problemática abordada, sino que también aporta evidencia para definir
las características funcionales y de seguridad de la solución
desarrollada.
